\section*{第 \ref{ch:b1:transient-regimes} 章——暂态：一阶与二阶}

\begin{solution}{exo:b1:transient-regimes:1}
$\tau = RC = 10^4 \times 47\times10^{-9} = \qty{0.47}{ms}$；$u_C(\tau)/E
= 1 - \eu^{-1} = 0.63$；$99\%$：$\eu^{-t/\tau} = 0.01$，$t = \tau\ln 100
= 4.6\tau = \qty{2.2}{ms}$。
\end{solution}

\begin{solution}{exo:b1:transient-regimes:2}
$\tau = L/R = \qty{4.0}{ms}$；$i_\infty = E/R = \qty{2.0}{A}$；$i(\tau) =
0.63 \times 2.0 = \qty{1.26}{A}$；在 $0^+$ 时，$i = 0$（连续性），因此
$u_L = E - Ri = \qty{10}{V}$。
\end{solution}

\begin{solution}{exo:b1:transient-regimes:3}
$\omega_0 = 1/\sqrt{LC} = \qty{3.16e4}{rad/s}$，$f_0 = \qty{5.03}{kHz}$；
$Q = \sqrt{L/C}/R = 316/100 = 3.2 > \tfrac12$：准周期的；
$\tau = 2Q/\omega_0 = \qty{0.20}{ms}$；$\omega = \omega_0\sqrt{1 - 1/4Q^2}
= 0.987\omega_0$，$T = \qty{0.20}{ms}$。
\end{solution}

\begin{solution}{exo:b1:transient-regimes:4}
$R_c = 2\sqrt{L/C} = \qty{632}{\ohm}$，$Q = \tfrac12$。当 \qty{2}{k\ohm} 时，
$Q = 0.16$：非周期的，缓慢衰减由较小的根
$\abs{r_+} \approx \omega_0 Q = \qty{5e3}{s^{-1}}$ 决定（时间尺度
\qty{0.2}{ms}，但不振荡）。
\end{solution}

\begin{solution}{exo:b1:transient-regimes:5}
$u = U_0\eu^{-t/\tau}$ 减半时，$\eu^{-t/\tau} = \tfrac12$：$t_{1/2}
= \tau\ln 2$。$\tau = 3.5/0.693 = \qty{5.05}{ms}$，$R = \tau/C =
\qty{5.1}{k\ohm}$。
\end{solution}

\begin{solution}{exo:b1:transient-regimes:6}
$i = (E/R)\eu^{-t/\tau}$：$\int_0^\infty Ri^2\,\dd t = (E^2/R)(\tau/2) =
\tfrac12 CE^2$，与 $R$ 无关。电源提供的能量为 $CE^2$：对任何电阻，效率均为
$50\%$。更好的办法是通过电感器充电（此时
电流不再与电压差成正比），或分若干
电压台阶——每步 $E/n$ 仅损失 $\tfrac12 C(E/n)^2$。
\end{solution}

\begin{solution}{exo:b1:transient-regimes:7}
$\delta = \ln(2.0/1.5) = 0.29$（且 $\ln(1.5/1.12) = 0.29$：一致）；
$Q \approx \pi/\delta = 11$。$T = \qty{0.40}{ms}$，故 $\omega_0 \approx
2\pi/T = \qty{1.57e4}{rad/s}$（修正项 $1/8Q^2$ 约为 $0.1\%$）。
$L = 1/(\omega_0^2 C) = \qty{4.1}{mH}$；$R = \omega_0 L/Q = \qty{5.8}{\ohm}$。
\end{solution}

\begin{solution}{exo:b1:transient-regimes:8}
$u_C \approx U_0\eu^{-t/\tau}\cos\omega_0 t$；$i = -C\,\dd u_C/\dd t
\approx CU_0\omega_0\sin\omega_0 t$（对 $t \ll \tau$），峰值为 $U_0\sqrt{C/L}
= 100 \times 0.1 = \qty{10}{A}$，在闭合后四分之一周期时达到，
此时 $u_C = 0$：全部能量（$\tfrac12 CU_0^2 = \qty{50}{mJ}$）此时为磁能，
$\tfrac12 Li^2 = \qty{50}{mJ}$。
\end{solution}

\begin{solution}{exo:b1:transient-regimes:9}
$\omega_0 = \sqrt{k/m} = \qty{7.1}{rad/s}$，$f_0 = \qty{1.1}{Hz}$；
$\alpha_c = 2\sqrt{km} = 2\sqrt{8\times10^6} = \qty{5.7e3}{N.s/m}$。
$\alpha = 2000$：$Q = \sqrt{km}/\alpha = 2828/2000 = 1.4$，准周期的；
$T = 2\pi/(\omega_0\sqrt{1 - 1/4Q^2}) = 0.89/0.935 = \qty{0.95}{s}$；
$\tau = 2Q/\omega_0 = \qty{0.40}{s}$。
\end{solution}

\begin{solution}{exo:b1:transient-regimes:10}
$i = I_0\eu^{-t/\tau}$，$\tau = L/R = \qty{5.0}{ms}$；峰值电压
为 $RI_0 = \qty{200}{V}$（$t = 0^+$ 时）；能量 $\tfrac12 LI_0^2 = \qty{1.0}{J}$
消耗在 $R$ 中。若无此电阻，电流便无通路：
$\dd i/\dd t$ 极大，线圈电压达到千伏量级，
并在开关处产生电弧。
\end{solution}

\begin{solution}{exo:b1:transient-regimes:11}
从 $C$ 看去：$E_{\mathrm{Th}} = ER_2/(R_1 + R_2) = \qty{8.0}{V}$，
$R_{\mathrm{Th}} = R_1 \parallel R_2 = \qty{6.67}{k\ohm}$；$\tau =
R_{\mathrm{Th}}C = \qty{6.7}{ms}$；$u_C = 8.0(1 - \eu^{-t/\tau})$ V。
\end{solution}

\begin{solution}{exo:b1:transient-regimes:12}
二重根 $-\omega_0$：$x = x_\infty + (A + Bt)\eu^{-\omega_0 t}$；
由 $x(0) = 0$ 得 $A = -x_\infty$，由 $x'(0) = 0$ 得 $B = \omega_0 A$：
$x = x_\infty[1 - (1 + \omega_0 t)\eu^{-\omega_0 t}]$。$95\%$：
$(1 + s)\eu^{-s} = 0.05$，$s \approx 4.7$：$t_{95} = 4.7/\omega_0$。对于
$Q = 5$ 时，包络 $\eu^{-t/\tau}$ 中 $\tau = 10/\omega_0$，需要
$3\tau = 30/\omega_0$；对于 $Q = 0.1$，慢根为 $\abs{r_+} \approx
\omega_0 Q$，$3/\abs{r_+} = 30/\omega_0$。两者都慢六倍：
临界阻尼是无超调的最快回复方式，
这正是仪表、闭门器或悬挂装置所需要的。
\end{solution}

\begin{solution}{pb:b1:transient-regimes:1}
\textbf{1.} $E = L\,\dd i/\dd t + Ri + u_C$，$i = C\,\dd u_C/\dd t$：
$LC\,u_C'' + RC\,u_C' + u_C = E$。

\textbf{2.} $u_C'' + (\omega_0/Q)u_C' + \omega_0^2u_C = \omega_0^2E$，其中
$\omega_0 = 1/\sqrt{LC}$，$Q = \sqrt{L/C}/R$。

\textbf{3.} $\omega_0 = \qty{3.16e4}{rad/s}$，$f_0 = \qty{5.03}{kHz}$，
$R_c = 2\sqrt{L/C} = \qty{632}{\ohm}$。

\textbf{4.} $r^2 + (\omega_0/Q)r + \omega_0^2 = 0$，$\Delta = \omega_0^2(1/Q^2
- 4)$：$Q < \tfrac12$ 非周期的，$Q = \tfrac12$ 临界阻尼，$Q > \tfrac12$
准周期的。

\textbf{5.} $\tau = 2Q/\omega_0$，$\omega = \omega_0\sqrt{1 - 1/4Q^2}$，
$u_C = E + \eu^{-t/\tau}(A\cos\omega t + B\sin\omega t)$。

\textbf{6.} $Q = 3.16$，$\tau = \qty{0.20}{ms}$，$T = 2\pi/\omega =
\qty{0.20}{ms}$，约三个振荡周期。

\textbf{7.} $u_C(0^+) = U_0$（$u_C$ 连续：电流有限）；
$u_C'(0^+) = i(0^+)/C = 0$（$i_L$ 连续，且此前为零）。

\textbf{8.} $m\ddot x = -k(x + x_{\text{平衡}}) - \alpha\dot x + mg$，其中
$kx_{\text{平衡}} = mg$：$m\ddot x + \alpha\dot x + kx = 0$。

\textbf{9.} $\omega_0 = \sqrt{k/m}$，$Q = \sqrt{km}/\alpha$；$x \leftrightarrow
q$，$\dot x \leftrightarrow i$，$m \leftrightarrow L$，$\alpha \leftrightarrow
R$，$k \leftrightarrow 1/C$。

\textbf{10.} $\omega_0 = \sqrt{22000/350} = \qty{7.93}{rad/s}$，$f_0 =
\qty{1.26}{Hz}$；$\alpha_c = 2\sqrt{km} = \qty{5.55e3}{N.s/m}$。

\textbf{11.} $Q = 2775/4000 = 0.69$，$\xi = 0.72$：准周期的，
略小于临界阻尼。

\textbf{12.} $\tau = 2Q/\omega_0 = \qty{0.18}{s}$；$\omega = \omega_0\sqrt{1
- \xi^2} = \qty{5.5}{rad/s}$，$T = \qty{1.1}{s}$。

\textbf{13.} 超调量 $\exp(-\pi \times 0.72/0.69) = \exp(-3.3) = 3.8\%$：
越过平衡位置后仅有一次小幅回升，随即静止——十分舒适。

\textbf{14.} $R = \sqrt{L/C}/Q = 316/0.69 = \qty{455}{\ohm}$。

\textbf{15.} $Q = 2775/1200 = 2.3$，$\xi = 0.22$；$\omega = 0.976\omega_0$，
$T = \qty{0.81}{s}$；$\tau = 2Q/\omega_0 = \qty{0.58}{s}$。

\textbf{16.} $\eu^{-T/\tau} = \eu^{-1.39} = 0.25$：每次回弹都是
前一次的四分之一；$0.25^3 = 1.6\%$：可见的回弹共三次。

\textbf{17.} $x(0) = -x_0$，$\dot x(0) = 0$；能量 $\tfrac12 kx_0^2 =
0.5 \times 22000 \times 0.0025 = \qty{28}{J}$，以热量形式
在阻尼器的油中耗散。

\textbf{18.} $v \approx x_0\omega_0 = 0.05 \times 7.93 = \qty{0.40}{m/s}$；
阻尼力 $\alpha v \approx 1200 \times 0.40 = \qty{480}{N}$。

\textbf{19.} 当 $Q = \tfrac12$ 时，回归过程为 $(1 + \omega_0 t)\eu^{-\omega_0
t}$，它渐近地缓慢趋近平衡：最后几毫米需要很长时间。一个
略欠阻尼系统（$\xi \approx 0.7$）超调几个
百分点，但能更快地进入并保持在允许误差带之内。

\textbf{20.} $\exp(-\pi \times 0.22/0.976) = \exp(-0.70) = 50\%$，与
$3.8\%$ 相比：磨损的汽车会越过平衡位置回弹一半。

\textbf{21.} $\delta = \ln 4 = 1.39$（两次，结果一致）；$\xi =
\delta/\sqrt{4\pi^2 + \delta^2} = 1.39/6.43 = 0.215$，$Q = 1/2\xi = 2.3$：
即磨损后的数值。

\textbf{22.} $-\ln 0.038 = 3.27 = \pi\xi/\sqrt{1 - \xi^2}$，故 $\xi =
3.27/\sqrt{\pi^2 + 3.27^2} = 0.72$，$Q = 0.69$。

\textbf{23.} 能量 $\propto$ 振幅$^2$：每个
周期剩余 $0.25^2 = 6\%$，故每个准周期耗散 $94\%$。

\textbf{24.} $R = 316/2.3 = \qty{137}{\ohm}$。在变量 $\omega_0 t$
中，标准方程只含 $Q$：两个具有相同 $Q$ 的系统
描出的曲线相同。

\textbf{25.} 仅品质因数 $Q$（或 $\xi$）即可决定形状；设计者
通常以 $Q \approx 0.7$（$\xi \approx 0.7$）为目标；若可见
的回弹多于一次，则 $Q$ 远大于 $1$：阻尼器已磨损。
\end{solution}
